\subsection{Симплекс-метод: идея и вспомогательная задача}

\begin{definition}[симплекс-метод]
\textbf{Симплекс-метод} — метод, основанный на \emph{направленном} переборе крайних
точек так, чтобы значение целевой функции в следующей точке было не хуже, чем в
предыдущей. Является итерационной схемой.
\end{definition}

\begin{idea}
Разница с полным перебором из следствия к теореме~\ref{th:optimal-is-support}: вместо
того чтобы обойти все $\binom{n}{m}$ базисов, на каждом шаге меняется ровно один
столбец базиса, причём так, чтобы целевая функция не возрастала. Монотонность
гарантирует, что ни один базис не повторится (при отсутствии вырожденности), а
конечность числа базисов — что алгоритм остановится. Подробный разбор шага —
в методичках.
\end{idea}

Чтобы запустить итерации, нужна стартовая крайняя точка. Её даёт вспомогательная
задача.

\subsubsection*{Вспомогательная задача симплекс-метода (метод искусственного базиса)}

Для исходной задачи
\[
  \begin{cases}
    c\T[N]\, x[N] \to \min, \\
    A[M,N]\, x[N] = b[M], \\
    x[N] \ge 0
  \end{cases}
\]
строим \textbf{вспомогательную задачу}, добавляя искусственные переменные $y[M]$:
\[
  \begin{cases}
    \displaystyle\sum_{i \in M} y[i] \to \min, \\[6pt]
    A[M,N]\, x[N] + E[M,M]\, y[M] = b[M], \\
    x[N] \ge 0, \quad y[M] \ge 0 ,
  \end{cases}
\]
где $E[M,M]$ — единичная матрица.

У вспомогательной задачи целевая функция \textbf{ограничена снизу} (она
неотрицательна), а множество допустимых точек \textbf{непусто}. Для неё сразу
известно начальное приближение:
\[
  x[N] = 0, \qquad y[M] = b[M] .
\]

К этому решению вспомогательной задачи применяем симплекс-метод. Если в результате
окажется, что $y[M] = 0$, то соответствующий $x[N]$ — опорный вектор исходной задачи,
с которого и стартуют основные итерации.

\begin{idea}
Начальное приближение $y[M] = b[M]$ требует $b[M] \ge 0$. Это не ограничение
общности: строку с $b_i < 0$ домножают на $-1$. Столбцы $E[M,M]$ линейно независимы,
поэтому точка $(x, y) = (0, b)$ — опорный вектор вспомогательной задачи, и запуск
корректен.

Дальше возможны два исхода:
\begin{itemize}[nosep,left=0pt]
  \item минимум вспомогательной задачи равен нулю, то есть $y[M] = 0$ — тогда
        найденный $x[N]$ допустим в исходной задаче и является её опорным вектором;
  \item минимум строго положителен — тогда допустимое множество исходной задачи
        \emph{пусто}, и исходная задача неразрешима.
\end{itemize}
Так вспомогательная задача одновременно даёт старт и проверяет совместность.
\end{idea}
